\setcounter{equation}{75}
\section{Extensional solution. Symmetrical conditions.}\label{article346}

When the middle surface is a right circular cone of angle \(2 \gamma \),
we take \(\alpha \) to be the distance \(r\) of a point on the surface
from the vertex of the cone, and then we have

\[\alpha = r, \quad A=1, \quad B=r \sin \gamma, \quad \frac{1}{R_1}=0,
\frac{1}{R_2}=\frac{1}{r \tan \gamma}, \quad \frac{d B}{d \alpha}=\sin \gamma.\]

When \(n=0\), the extensional solutions include the solution for
torsion, which may be recorded here in the form
\begin{align}
s_1&=\frac{a_1}{r^2}-\frac{1-\sigma^2}{2 E h r^2} \int_{r_0}^r r^2 Y^{\prime} d r\\
V&=b_1 r+\frac{2 r}{1-\sigma} \int_{r_0}^r \frac{s_1}{r} d r \nonumber \label{eq76}
\end{align}
and a solution in which \(V\) and \(s_1\) vanish. In equations
(76) \(a_1\) and \(b_1\) are constants of integration, and \(r_0\) a
suitably chosen fixed value of \(r\).

For the extensional solution in which \(V\) and \(s_1\) vanish we have
the equations
\begin{equation}
\left.\begin{array}{r}
\frac{d}{d r}\left(t_1 r\right)-t_2+\frac{1-\sigma^2}{2 E h}+X^{\prime \prime}=0,\\
\frac{t_2}{r \tan \gamma}+\frac{1-\sigma^2}{2 E h} Z^{\prime \prime}=0,\label{eq77}
\end{array}\right\}
\end{equation}
with
\begin{align}
& t_1=\frac{d U}{d r}+\sigma\left(\frac{U}{r}-\frac{W}{r \tan \gamma}\right), \nonumber \\
& t_2=\sigma \frac{d U}{d r}+\left(\frac{U}{r}-\frac{W}{r \tan \gamma}\right).\label{eq78}
\end{align}
From equations~(\ref{eq77}) we find % eq77 = 82
\begin{align}
& t_1=\frac{a_3}{r}-\frac{1-\sigma^2}{2 E k r} \int_r^r r
\left(X^{\prime \prime}+ Z^{\prime \prime} \tan \gamma \right) d r, \nonumber \\
& t_2=-\frac{1-\sigma^2}{2 E h}+Z^{\prime \prime} \tan \gamma, \label{eq79}
\end{align}
where \(a_2\) is a constant of integration. Thus \(t_1\) and \(t_2\) are
known.  From equations~(\ref{eq78}) we then find  % eq78 = 83
\begin{align}
& U=b_t+\int_{r_t}^r \frac{t_1-\sigma t_2}{1-\sigma^2} d r_2 \nonumber \\
& W=U \tan \gamma-r \tan \gamma \frac{t_2-\sigma t_1}{1-\sigma^2} \label{eq80}
\end{align}
where \(b_2\) is a constant of integration. Thus \(U\) and \(W\) are known.
\section{Edge-effect. Symmetrical conditions.}\label{article345}
As in previous examples we take \(s_1, s_2\) and \(h_1\) to vanish, and
omit external forces. We have the equations
\begin{align}
\frac{u}{d r}\left(t_1 r\right)-t_2&=0, \nonumber \\
n_2&=0, \nonumber \\
\frac{d}{d r}\left(n_1 r\right)+\frac{t_2}{\tan \gamma}&=0,\label{eq81} \\
\frac{d}{d r}\left(g_1 r\right)-g_2-n_1 r&=0 \nonumber .
\end{align}
From the \nth{1} and \nth{3} of equations~(\ref{eq81}) we obtain the equation
\[t_1 r+n_1 r \tan(\gamma)=\text { const },\]
and, as before, we equate the constant in the right-hand member to zero,
thus obtaining the equation
\begin{equation}\label{eq82} n_1=-t_1 \cot \gamma_2 \end{equation}
Then there remain two equations, viz.
\begin{align}
& t_2=\frac{d}{d r}\left(t_1 r\right), \nonumber \\
& g_2=\frac{d}{d r}\left(g_1 r\right)+t_2 r \cot \gamma \label{eq83}
\end{align}
Again, we have the equations
\begin{align}
& t_1=\frac{d U}{d r}+\sigma\left(\frac{U}{r}-\frac{W}{r \tan \gamma}\right),\nonumber \\
& t_2=\sigma \frac{d U}{d r}+\left(\frac{U}{r}-\frac{W}{r \tan \gamma}\right),\nonumber \\
& g_1=-\frac{h^2}{3}\left(\frac{d^2 W}{d r^2}+\frac{\sigma}{r} \frac{d W}{d r}\right),\label{eq84} \\
& g_2=-\frac{h^2}{3}\left(\sigma \frac{d^2 W}{d r^2}+\frac{1}{r} \frac{d W}{d r}\right). \nonumber
\end{align}
In these we write \(\xi \) for \(d W / d r\). The \nth{2} of equations~(\ref{eq83})
becomes, on substituting for \(g_1, g_2\) from the \nth{3} and \nth{4} of
equations~(\ref{eq84}) and using~(\ref{eq82}),
\[
\frac{h^2}{3}\big(\frac{d}{d r}\left(r \frac{d \xi}{d r}+
\sigma \xi \right)-\left(\sigma \frac{d \xi}{d r}+
\frac{\xi}{r}\right)\big)+n_1 r=0,
\]
or
\begin{equation}\label{eq85}
\frac{h^2}{3}\left(\frac{d^2 \xi}{d r^2}+
\frac{1}{r} \frac{d \xi}{d r}-\frac{\xi}{r^2}\right)+n_1=0.
\end{equation}
Now the \nth{1} and \nth{2} of equations~(\ref{eq84}) give
\begin{align*}
U & =W \cot \gamma+r \frac{t_2-\sigma t_1}{1-\sigma^2}, \\
\frac{d U}{d r} & =\frac{t_1-\sigma t_2}{1-\sigma^2}.
\end{align*}

We eliminate \(U\) from these equations, thus obtaining the equation
\[
\left(1-\sigma^2\right) \xi \cot(\gamma) +
\frac{d}{d r}\left( r\left(t_1-\sigma t_1\right) \right)
-\left(t_1-\sigma t_2\right)=0,\]
which becomes, on substituting for \(t_2\) from the \nth{1} of
equations~(\ref{eq83}) and for \(t_1\) from~(\ref{eq82}),
\[
\frac{d}{d r}\left( r \frac{d}{d r} (n_1 r) ~~ \right)-
n_1-\left(1-\alpha^2\right) \xi \cot ^2 \gamma=0,
\]
or
\begin{equation}\label{eq86}
\frac{h^2}{3}\left(\frac{d^2 \xi}{d r^2}+
\frac{1}{r} \frac{d \xi}{d r}-\frac{\xi}{r^2}\right)+n_1=0.
\end{equation}
On putting
\begin{equation}\label{eq87}
n_1 r = \eta, \quad r = x^2/2,
\end{equation}
equations~(\ref{eq85}) and~(\ref{eq86}) become
\begin{align}
& \frac{d^2 \xi}{d x^2}+\frac{1}{x} \frac{d \xi}{d x}-\frac{4}{x^2} \xi+\frac{6}{h^2} \eta=0, \nonumber \\
& \frac{d^2 \eta}{d x^2}+\frac{1}{x} \frac{d \eta}{d x}-\frac{4}{x^2} \eta-\frac{2\left(1-\sigma^2\right)}{\tan^2(\gamma)} \xi=0 \label{eq88}.
\end{align}
To solve these equations put
\[
-\;\frac{2\,(1-\sigma^{2})}{\tan^{2}\gamma}\; \xi=k^{2}\eta, \qquad
\frac{6}{h^{2}}\,\eta=k^{2}\xi,
\]
so that
\begin{equation}\label{eq89}
k^{4}=-\,{\frac{(1-\sigma^{2})}{\hbar^{4}\tan^{3}\gamma}}\,.
\end{equation}
then the integrals of the equations are Bessel's functions of
\(k x\)
of order 2.
There are four complex values of
\(k\), but two of them are irrelevant, as the Bessel's functions are not altered by changing the sign of
\(k\).
The two relevant values of
\(k\)
are
\[
k= (1\pm i){(3(1-\sigma^{2}))}^{\frac{1}{4}}/{(h\tan \gamma )}^{\frac{1}{2}},
\]
or say
\[
k=\left(1\pm i\right)m^{\frac{1}{2}},
\]
where
\begin{equation}
m= \frac{\sqrt{ 3\left(1-\sigma^{2}\right)}}{h\tan\gamma}
\end{equation}
\[
\]
The integrals of the equations~(\ref{eq88}) are then
\begin{align}
\xi=&A_{1}J_{2}\left((1+i){(2mr)}^{\frac{1}{2}}\right)+
B_{1}J_{2}\left((1+i)\ {(2 mr)}^{\frac{1}{2}}\right) \nonumber \\
&+A_{2}J_{2}\left((1-i)\ {(2mr)}_{*}^{\frac{1}{2}}\right)+
B_{2}Y_{2}\left((1-i){(2mr)}^{\frac{1}{2}}\right) \nonumber \\
\eta&=\frac{1-\sigma^{2}}{m~\tan^{2}\gamma}\,i
\left( A_{1}J_{2}\left((1+i)\,{(2mr)}^{\frac{1}{2}}\right)+
B_{1}Y_{2}\,\left((1+i)\,{(2mr)}^{\frac{1}{2}}\right)\right) \label{eq91}\\
&-\frac{1-\sigma^{2}}{m~\tan^{2}\gamma}\,i
\left( A_{2}J_{2}\, \{(1-i)\,{(2mr)}^{\frac{1}{2}}\}
+B_{2}Y_{2}\, \{(1-i)\,{(2mr)}^{\frac{1}{2}}\} \right) \nonumber .
\end{align}

In equations~(\ref{eq91}) the constants \(A_i\), \(B_i\) are complex and may be
adjusted so that the values of
\(\xi \) and \(\eta \) are
real. By means of the semiconvergent expansions of the Bessel's functions we
can see that the edge-effect diminishes with increasing distance from an edge
in much the same way as it has been found to do in previous examples.


The integration of the equations for a conical shell under symmetrical forces is due to E. Meissner*,
who has also pointed out that for a shell of variable thickness, proportional
to distance from the vertex of the cone, the integration can be effected
without introducing any transcendental functions. A detailed investigation for
the case of uniform thickness and for a number of laws of variable thickness
has been given by F. Dubois \(\dagger \).

When
\(\xi \) and \(\eta \) are known,
\(d W/dr\) and \(n_1\) are known, and by equation (82) and the 1st of equations (83)
\(t_1\) and \(t_2\) are known.  Also
\(g_1\) and \(g_2\) are known, and thus all the stress-resultants and
stress-couples are known.  To determine the displacement an additional
integration is required, in order to find
\(W\) from the known value of
\(\xi \) or \(d W/dr\).  When
\(W\) is known,
\(U\) is known from the equation giving
\(t_2 - \sigma t_1\).  In effecting the integration for
\(W\) no constant need be added, as an additive constant
be added, as an additive constant in the expression for
\(W\) occurs also in the extensional solution.

* loc.\ cit.\ ante, p. 581.

\(\dagger \) ``Über die Festigkeit der Kegelschale'' (Diss.), Zürich, 1917.
